\section{Kidney Yogurt}
Over time, our microbiomes have been modified by external forces. These include non-standard diets, antibiotics and other factors which have eliminated some microbes and allowed so-called "bad bacteria" to thrive. This is a burgeoning field and is likely to change dramatically as we learn more about what is effective. One approach, pioneered by Dr. William Davis, is to ferment yogurts and taylor-make them to contain large colonies of the desired species. This is actually a very simple process though it must be done properly to ensure success. The objective is to select bacterial species and strains to obtain specific health outcomes. Commercial yogurts are not designed to do this so the approach is a DIY one. The fermentation process is modified to  exponentially increase bacterial counts. The logic here is more is better. The results of the yogurt on the targeted health objective will determine whether that logic is correct. With these fermentation projects, we can get bacterial counts in the hundreds of billions. Commercial yogurts are typically in the millions or at a few billion at the highest. High counts are achieved by extending the fermentation time from four hours for commercial yogurts to typically 36 hours. This results in a thousand fold difference. It is vital to add a prebiotic fiber to provide the nourishment to the microbes. That also helps increase the richness and thickness of the final product.

One benefit to the prolonged fermentation times is it reduces lactose, the milk sugar in dairy, to negligible levels as it is converted to lactic acid, which can lend a tart flavor to a yogurt fermented in this wise.
